% GENERATED FILE. Edit canonical lessons, then run npm run generate:books.
% canonical-source-hash: fnv1a64:1afe9762819c2507
% canonical-lessons: MW-C136-mulo, MW-C136-lasan, MW-C136-adrak, MW-C136-haldi, MW-C136-jiro

\chapter{Radish, Garlic, Ginger, Turmeric, Cumin}
\label{ch:mw-radish-garlic-ginger-turmeric-cumin}
\hlchaptermodality{136}

\begin{chapteropening}
\textbf{By the end of this chapter:} \emph{I can say a radish, garlic, ginger, turmeric and cumin.}

Complete the last lesson of chapter 136: I can say a radish, garlic, ginger, turmeric and cumin.
\end{chapteropening}

\section[\texorpdfstring{\mw{मूळो} (mūḷo)}{मूळो (mūḷo)}]{\mw{मूळो} (mūḷo) — a radish}

\label{lesson:MW-C136-mulo}



\begin{quote}
Before the new one: say the Marwadi for an aubergine, then the Marwadi for spinach.
\end{quote}

\subsection*{You'll want to know: \mw{मूळो}}
\textbf{\mw{मूळो}} — \emph{mūḷo} — \textquotedblleft{}a radish\textquotedblright{}.

\begin{cousinweb}
\textbf{\mw{मूळो}} (\emph{mūḷo}) is a tadbhava, a word inherited from Sanskrit and reshaped by sound change. It comes from \emph{mūla}, \textquotedblleft{}root; radish\textquotedblright{}. CDIAL lists Old Marwari \emph{mūlī} under it (CDIAL 10250). CDIAL notes a likely non-Aryan origin for the word; Gujarati has \emph{mūḷo}, \textquotedblleft{}radish\textquotedblright{}.
\end{cousinweb}

\begin{grammarlens}[title={What you've built}]
One more word for this chapter.
\end{grammarlens}

\subsection*{Your turn}
\begin{itemize}
\raggedright
  \item \emph{Say it:} \emph{mūḷo}
  \item \emph{Say it:} \emph{mūḷo}, once more
  \item \emph{Say it:} say \emph{mūḷo}
  \item \emph{Recall:} say the Marwadi for an aubergine, then the Marwadi for spinach, then say \emph{mūḷo} again
\end{itemize}

\begin{tcolorbox}[breakable,colback=teal!4,colframe=teal!35!black,title={Before you move on}]
What does \mw{मूळो} mean? (\textquotedblleft{}A radish\textquotedblright{}.) Say it once more.
\end{tcolorbox}
\section[\texorpdfstring{\mw{लसण} (lasaṇ)}{लसण (lasaṇ)}]{\mw{लसण} (lasaṇ) — garlic}

\label{lesson:MW-C136-lasan}



\begin{quote}
Before the new one: say the Marwadi for spinach, then the Marwadi for a radish.
\end{quote}

\subsection*{You'll want to know: \mw{लसण}}
\textbf{\mw{लसण}} — \emph{lasaṇ} — \textquotedblleft{}garlic\textquotedblright{}.

\begin{cousinweb}
\textbf{\mw{लसण}} (\emph{lasaṇ}) is a tadbhava, a word inherited from Sanskrit and reshaped by sound change. It comes from \emph{laśuna}, \textquotedblleft{}garlic\textquotedblright{}. CDIAL gives no Marwari form in this entry (CDIAL 10990); the match is through Gujarati \emph{lasaṇ}.
\end{cousinweb}

\begin{grammarlens}[title={What you've built}]
One more word for this chapter.
\end{grammarlens}

\subsection*{Your turn}
\begin{itemize}
\raggedright
  \item \emph{Say it:} \emph{lasaṇ}
  \item \emph{Say it:} \emph{lasaṇ}, once more
  \item \emph{Say it:} say \emph{lasaṇ}
  \item \emph{Recall:} say the Marwadi for spinach, then the Marwadi for a radish, then say \emph{lasaṇ} again
\end{itemize}

\begin{tcolorbox}[breakable,colback=teal!4,colframe=teal!35!black,title={Before you move on}]
What does \mw{लसण} mean? (\textquotedblleft{}Garlic\textquotedblright{}.) Say it once more.
\end{tcolorbox}
\section[\texorpdfstring{\mw{अदरक} (adrak)}{अदरक (adrak)}]{\mw{अदरक} (adrak) — ginger}

\label{lesson:MW-C136-adrak}



\begin{quote}
Before the new one: say the Marwadi for a radish, then the Marwadi for garlic.
\end{quote}

\subsection*{You'll want to know: \mw{अदरक}}
\textbf{\mw{अदरक}} — \emph{adrak} — \textquotedblleft{}ginger\textquotedblright{}.

\begin{cousinweb}
\textbf{\mw{अदरक}} (\emph{adrak}) goes back to Sanskrit \emph{ārdraka}, \textquotedblleft{}fresh ginger\textquotedblright{} (CDIAL 1341). CDIAL lists the fully inherited forms as Hindi \emph{ādā} and Gujarati \emph{ādũ}; \textbf{\mw{अदरक}} (\emph{adrak}) keeps the \emph{r} and \emph{k} those forms lost, and CDIAL does not list it.
\end{cousinweb}

\begin{grammarlens}[title={What you've built}]
One more word for this chapter.
\end{grammarlens}

\subsection*{Your turn}
\begin{itemize}
\raggedright
  \item \emph{Say it:} \emph{adrak}
  \item \emph{Say it:} \emph{adrak}, once more
  \item \emph{Say it:} say \emph{adrak}
  \item \emph{Recall:} say the Marwadi for a radish, then the Marwadi for garlic, then say \emph{adrak} again
\end{itemize}

\begin{tcolorbox}[breakable,colback=teal!4,colframe=teal!35!black,title={Before you move on}]
What does \mw{अदरक} mean? (\textquotedblleft{}Ginger\textquotedblright{}.) Say it once more.
\end{tcolorbox}
\section[\texorpdfstring{\mw{हळदी} (haḷdī)}{हळदी (haḷdī)}]{\mw{हळदी} (haḷdī) — turmeric}

\label{lesson:MW-C136-haldi}



\begin{quote}
Before the new one: say the Marwadi for garlic, then the Marwadi for ginger.
\end{quote}

\subsection*{You'll want to know: \mw{हळदी}}
\textbf{\mw{हळदी}} — \emph{haḷdī} — \textquotedblleft{}turmeric\textquotedblright{}.

\begin{cousinweb}
\textbf{\mw{हळदी}} (\emph{haḷdī}) is a tadbhava, a word inherited from Sanskrit and reshaped by sound change. It comes from \emph{haridrā}, \textquotedblleft{}turmeric\textquotedblright{}. CDIAL gives no Marwari form in this entry (CDIAL 13992); the match is through Hindi \emph{haldī}. Marwadi says the \emph{l} curled back.
\end{cousinweb}

\begin{grammarlens}[title={What you've built}]
One more word for this chapter.
\end{grammarlens}

\subsection*{Your turn}
\begin{itemize}
\raggedright
  \item \emph{Say it:} \emph{haḷdī}
  \item \emph{Say it:} \emph{haḷdī}, once more
  \item \emph{Say it:} say \emph{haḷdī}
  \item \emph{Recall:} say the Marwadi for garlic, then the Marwadi for ginger, then say \emph{haḷdī} again
\end{itemize}

\begin{tcolorbox}[breakable,colback=teal!4,colframe=teal!35!black,title={Before you move on}]
What does \mw{हळदी} mean? (\textquotedblleft{}Turmeric\textquotedblright{}.) Say it once more.
\end{tcolorbox}
\section[\texorpdfstring{\mw{जीरो} (jīro)}{जीरो (jīro)}]{\mw{जीरो} (jīro) — cumin}

\label{lesson:MW-C136-jiro}



\begin{quote}
Before the new one: say the Marwadi for ginger, then the Marwadi for turmeric.
\end{quote}

\subsection*{You'll want to know: \mw{जीरो}}
\textbf{\mw{जीरो}} — \emph{jīro} — \textquotedblleft{}cumin\textquotedblright{}.

\begin{cousinweb}
\textbf{\mw{जीरो}} (\emph{jīro}) is a tadbhava, a word inherited from Sanskrit and reshaped by sound change. It comes from \emph{jīraka}, \textquotedblleft{}cummin-seed\textquotedblright{}. CDIAL gives no Marwari form in this entry (CDIAL 5234); the match is through Hindi \emph{jīrā} and Gujarati \emph{jīrũ}.
\end{cousinweb}

\begin{grammarlens}[title={What you've built}]
One more word for this chapter.
\end{grammarlens}

\subsection*{Your turn}
\begin{itemize}
\raggedright
  \item \emph{Say it:} \emph{jīro}
  \item \emph{Say it:} \emph{jīro}, once more
  \item \emph{Say it:} say \emph{jīro}
  \item \emph{Recall:} say the Marwadi for ginger, then the Marwadi for turmeric, then say \emph{jīro} again
\end{itemize}

\begin{tcolorbox}[breakable,colback=teal!4,colframe=teal!35!black,title={Before you move on}]
What does \mw{जीरो} mean? (\textquotedblleft{}Cumin\textquotedblright{}.) Say it once more.
\end{tcolorbox}
